% SYNTHETIC TEST FIXTURE for paper-copilot.
% Built from three artifacts of powerofjinbo/Displaced_tracks_GNN (results/addendum_g_g1_v1_claim.json,
% results/addendum_g_g1_v1.csv, docs/definitions.md). It is not the Canto manuscript; it exercises
% segmentation, grounding, linting and compilation.
\documentclass[aps,prd,twocolumn,superscriptaddress,nofootinbib]{revtex4-2}
\usepackage{amsmath}
\usepackage{graphicx}
\usepackage{hyperref}
\newcommand{\ttbar}{\ensuremath{t\bar{t}}}

\begin{document}

\title{Fixture: tracking displaced particles with a graph neural network}
\author{A. Fixture}
\affiliation{Nowhere Institute}

\begin{abstract}
Charged particles from the decay of a long-lived particle start far from the collision point.
Trackers tuned on prompt particles lose many of them.
However, graph-neural-network trackers have not been trained without the signal they are meant to find.
We study a tracker trained only on \ttbar{} events.
On isolated displaced particles that neither model saw in training, it finds 96.94\% of the targets, compared with 94.62\% for the signal-trained baseline, at a fake-track rate of 0.55\% against 1.38\%.
\end{abstract}

\maketitle

\section{Introduction}
\label{sec:intro}
Searches for long-lived particles look for tracks that do not point back to the beam line~\cite{Ju2022Zenodo}.
Machine-learning pipelines such as Exa.TrkX were built for prompt tracks, and their large-radius extension was studied by Wang et al.~\cite{Wang2022LRT}.
However, these studies trained on the signal they evaluate.
In this paper, we train on generic events only and test on particles generated after the design freeze.
Section~\ref{sec:method} describes the selection, and Sec.~\ref{sec:results} shows the result.

\section{Method}
\label{sec:method}
A target particle has $p_T > 1$~GeV, $|\eta| < 4$ and at least five space points.
A reconstructed track matches a particle if more than half of its hits come from that particle and more than half of the particle's hits are on the track, i.e., double majority.
The efficiency is
\begin{equation}
  \epsilon = \frac{N_{\text{matched}}}{N_{\text{targets}}},
  \label{eq:eff}
\end{equation}
where both counts use the same population.
The uncertainty in Eq.~(\ref{eq:eff}) combines the event bootstrap and the spread of three training seeds.\footnote{The bootstrap uses one fixed list of resampled events per test set.}

\section{Results}
\label{sec:results}
On the isolated gun particles, the tracker finds 2.32~$\pm$~0.28 points more than acorn (HNL) with the same track builder; see Fig.~\ref{fig:eff}.
The test set holds 8,710 targets.
The fake-track rate is 0.55\% against 1.38\% for the baseline.
This lead holds for every particle species.

\begin{figure}[t]
  \centering
  \fbox{\parbox{0.8\linewidth}{\centering placeholder artwork}}
  \caption{Efficiency versus transverse displacement for the isolated gun particles. Solid: the \ttbar-trained tracker; dashed: acorn (HNL). A ratio above one means the new tracker finds more targets.}
  \label{fig:eff}
\end{figure}

\begin{table}[b]
  \caption{Efficiency on the isolated gun particles, in percent.}
  \label{tab:eff}
  \begin{tabular}{lc}
    Model & Efficiency \\
    tracker (\ttbar) & 96.94 \\
    acorn (HNL) & 94.62 \\
  \end{tabular}
\end{table}

\section{Conclusions}
\label{sec:concl}
We have shown that a tracker trained without the signal finds more displaced particles than one trained on it.
The principal limitation is that all results are within one simulated detector without pileup.

\appendix
\section{Validation}
\label{app:val}
% A comment with a period. It must not create a sentence.
On validation, the same comparison gave a smaller difference (see Table~\ref{tab:eff}).

\bibliography{refs}

\end{document}
